\subsection{Actual torsion and cusp corrections in the remaining projective self blocks}

\begin{statement}{D.62}{Proved}{Two source-defined correction maps lift all twenty-six surviving projective self-wedge classes}
Retain the actual sources, integral projectors and
homology action \(T_{61}\) of D.61. In particular
\[
 X=\widehat B/(R_{50}\cap L_2),\qquad
 Y=X/\mathcal K,\qquad
 X_{61}=X/\mathcal J_{61}.
 \tag{D.62.1}
\]
The following projector uses the one excluded
old self block and all eight new self blocks:
\[
 \begin{aligned}
 \Pi_{62}&=e_\Delta\left(
       \bigwedge\nolimits^2\bar e_{1,4}
          +\sum_{b=1}^{2}\sum_{i=1}^{4}
                    \bigwedge\nolimits^2\bar f_{b,i}\right),\\
 \mathcal C_{62}&=\Pi_{62}Z,\qquad
 \mathcal J_{62}=(1-\Pi_{62})\mathcal J_{61}.
 \end{aligned}
 \tag{D.62.2}
\]
These are actual equivariant central submodules,
and
\[
 \begin{gathered}
 \mathcal J_{61}=\mathcal C_{62}\oplus\mathcal J_{62},
 \qquad\dim\mathcal C_{62}=6656,\qquad
 \dim(\mathcal C_{62})_G=28,\\
 \operatorname{rank}\nu|_{(\mathcal C_{62})_G}=2,
 \qquad\dim\ker\nu|_{(\mathcal C_{62})_G}=26.
 \end{gathered}
 \tag{D.62.3}
\]

There is an actual continuous peripheral
automorphism of
\[
                         X_{62}=X/\mathcal J_{62}.
 \tag{D.62.4}
\]
It has multiplier one, fixes the original cusp
\(t\), induces identity on the arithmetic group
\(G\), and retains the complete D.57 action
and its literal compatible conjugators.
Its prescribed \(\mathcal W_h\) action,
normalized parameter, and all established
metabelian and CM comparisons through that
map are unchanged. Its integral homology
action is exactly \(T_{61}\).

The comparison on the full homology source
is explicit. For actual source elements
\(u_j\in S_{57}\), set
\[
 b_{t,j}=\prod_{k=0}^{35}t^k u_jt^{-k},
 \qquad b_t=\prod_j b_{t,j}.
 \tag{D.62.5}
\]
The homology vectors and the finite choice
of factors are specified in the proof.
The new \(Y\)-action is
\[
                 \operatorname{Int}(b_{t,Y})\alpha_{61,Y}.
 \tag{D.62.6}
\]
Here \(b_{t,Y}\) centralizes \(t\),
has zero image in \(X_*\) and
\(\mathcal W_h\), and belongs to the
abelian kernel of \(Y\to X_*\).
The other corrections leave the action
on \(Y\) unchanged. Thus (D.62.6)
gives the exact change of representative;
no unchanged extension of the chosen
D.61 representative is asserted.

The two actual correction maps on a
projective self factor \(E\) are
\[
 \begin{aligned}
 \eta_s(v)&=[e_\Delta(v\wedge sv)]_G,\\
 \eta_t(v)&=[e_\Delta(v\wedge N_{36}(t)v)]_G,
       &N_{36}(t)&=\sum_{k=0}^{35}t^k.
 \end{aligned}
 \tag{D.62.7}
\]
They are linear maps over \(\mathbb F_2\),
realized by the source words in the proof.
Their joint image is precisely the
twenty-six-dimensional kernel in (D.62.3).
The ranks on the three self types are
\[
 \begin{array}{c|r|r|r|r|r}
 \text{self type}&\text{copies}&\dim M_G
   &\operatorname{rank}\nu&\operatorname{rank}\eta_s
   &\operatorname{rank}\eta_t\\\hline
 (0,0,0,1),\ i=4&1&4&2&0&2\\
 (0,0,1,1)&4&3&0&2&1\\
 (0,0,1,2)&4&3&0&2&1
 \end{array}
 \tag{D.62.8}
\]
In each row the two correction images
have the full joint rank
\(\dim M_G-\operatorname{rank}\nu\).

The cusp map uses the literal affine
element in the actual order-46656
quotient:
\[
 \begin{aligned}
 t&=((1,1,0,1),q_4),\\
 t^{18}&=((0,2,1,0),q_{314}).
 \end{aligned}
 \tag{D.62.9}
\]
The matrices of \(q_4,q_{314}\)
are those of the retained marked state
list. The word \(t^{18}\) in the
full source is used with its actual
square \(t^{36}\); it is never
replaced by an order-two source element.

The exact source kernels are
\[
 \begin{aligned}
 \ker(X_{62}\to X_{61})&=\mathcal C_{62},\\
 \ker(X_{62}\to Y)&=(1-e_\Delta)\mathcal K
    \oplus\mathcal C_{59}\oplus\mathcal C_{60}
    \oplus\mathcal C_{61}\oplus\mathcal C_{62},\\
 \ker(X\to X_{62})&=\mathcal J_{62},\\
 \dim\mathcal J_{62}&=48w-2560-5308416\kappa-40960.
 \end{aligned}
 \tag{D.62.10}
\]
The specified remaining small-character
submodule has dimension 325120 and
coinvariant dimension 534. In the full
source put
\[
 S_{62}=\{b\in R_{50}:
        \operatorname{im}_{\mathcal T_2}(b)\in\mathcal J_{62}\}.
 \qquad
 S_{62}/(R_{50}\cap L_2)=\mathcal J_{62}.
 \tag{D.62.11}
\]
This remaining kernel is nonzero.
Its arithmetic boundary and genuine
full profinite realization remain
unresolved. The exponent-eight word
check below verifies the correction
identities, without asserting an
arithmetic lift to that larger quotient.
The manuscript passages are Part D(e),
section 51.16, folio 734 verso, scan 174,
and section 51.15, folios 725--727,
scans 156--160.
\end{statement}
\begin{proof}
\emph{The actual integral source of every correcting vector.}
Let \(E\) be one of the nine integral
projective summands appearing in (D.62.2),
and let \(e\) be its integral projector.
D.60 and D.61 prove
\[
 eT_{61}=e,\qquad e\mathcal R=0,
 \qquad\theta_U E=0,\qquad E\subset U^-.
 \tag{D.62.12}
\]
Here \(\mathcal R\) is the full cusp
homology lattice. D.58 proves the actual
surjection
\[
 S_{57}=\mathfrak R_h\cap M_*\cap L_{*,2}
       \longrightarrow
 \mathcal L=\ker(U^-\to U_*^-).
 \tag{D.62.13}
\]
For every vector of \(\bar E=E/2E\),
choose an integral lift in \(E\)
and an actual preimage \(u\in S_{57}\)
under (D.62.13). Thus \(u\) has zero
image in the complete \(X_*\) and
\(\mathcal W_h\) sources, with precisely
the prescribed full homology vector.
Its image in \(Y\) is the corresponding
element of the abelian kernel
\(\mathcal L=\ker(Y\to X_*)\).
The use of this surjection retains the
entire arithmetic source, rather than
choosing only a vector in a finite
coefficient model.

Lift the two conjugators of the chosen
\(Y\)-action to \(X\). Their torsion
conjugates define an actual homomorphism
\(\varphi:\widehat B\to X\), with
\[
 \widetilde s=\varphi(s),\quad
 \widetilde r=\varphi(r),\quad
 \widetilde s^2=\widetilde r^3=1,
 \qquad\varphi(t)=zt,\quad z\in\mathcal K.
 \tag{D.62.14}
\]
This source map induces the given
\(Y\)-action and its homology
\(T_{61}\). Its action on the central
module is \(\bigwedge^2\bar T_{61}\).

We use throughout the convention
\([a,b]=aba^{-1}b^{-1}\).
Commutators of the correcting kernel
elements below lie in the actual
central module \(Z\). Indeed their
\(\mathcal W_h\) image is one and
their full homology image is zero;
D.54 identifies this kernel in \(X\)
with \(Z\), faithfully through
\(\mathcal T_2\). Within that target
their coordinates are precisely the
exterior products of their homology
vectors. This justifies every following
class-two word identity inside \(X\).

\emph{An actual source-involution correction.}
For the chosen preimage \(u\), put
\[
 b_s=u\,\widetilde s u\widetilde s^{-1}.
 \tag{D.62.15}
\]
One may form this word in the full source:
choose a source lift \(\gamma\) of the
conjugator defining \(\widetilde s\) and
use the actual torsion element
\(\gamma s\gamma^{-1}\).
Normality of \(S_{57}\) then places
the entire correcting word in \(S_{57}\).
Its image in \(Y\) is \(v+sv\),
which is fixed by \(s\). Thus it
centralizes \(\widetilde s_Y\) in
\(Y\). Replacing \(\widetilde s\)
by \(b_s\widetilde s b_s^{-1}\)
leaves the \(Y\)-action unchanged.
It also retains the literal conjugator
in \(X_*\) and \(\mathcal W_h\),
since \(b_s\) has zero image there.
The exact identity \(\widetilde s^2=1\)
gives
\[
 b_s\widetilde s b_s^{-1}\widetilde s^{-1}
   =[u,\widetilde s u\widetilde s^{-1}].
 \tag{D.62.16}
\]
Its exterior coordinate is \(v\wedge sv\).
Consequently this change adds precisely
\(\eta_s(v)\) to the projected
coinvariant cusp error.

The map on coinvariants is linear.
The polar expression is
\(v\wedge sw+w\wedge sv\).
Applying \(s\) to its first term
gives the second in characteristic
two, so their classes cancel.
The map factors through \(N_2(s)v\):
for \(N_2(s)v=0\) one has \(sv=v\),
so its value is zero. Its full norm
is zero, since pairing \(g\) and
\(gs\) in the group sum cancels
their two exterior terms. These
statements remain true after the
equivariant \(\Delta\)-projection.

The first complete test,
\path{d_self_wedge_torsion_correction.py},
computes this actual source-\(s\) map
on all global basis vectors and checks
every polar equation. Its ranks are
the \(\eta_s\) column of (D.62.8).
It also preserves, as a separate
diagnostic, the maps for the other
three involution classes of \(Q\).
The diagnostic class 314 is not
silently used as another source
torsion generator. Under the actual
map \(\widehat B\to Q\), every
order-two source element is conjugate
to \(s\), by
\cite[Theorem A$'$]{A_HerfortRibes1985}.
The marked enumeration gives class
sizes 36 for \(s\) and three for
state 314; they are not conjugate.
Hence there is no order-two source
element projecting to that state.

The intermediate calculation
\path{d_self_wedge_coordinate_blocks.py}
records invariant coordinate summands
of each new self module. Their global
dimensions are \(96,96,288,288\).
The diagonal 96-dimensional summand
and the \((0,5)\)-orbit summand of
dimension 288 have onto source-\(s\)
correction on their two coinvariants;
one further 96-dimensional summand
has zero coinvariants. This proves
the exact sixteen-dimensional image
of the first route while preserving
its then-uncontrolled ten directions.
The cusp construction supplies these
directions using a different actual
source word.

\emph{The actual infinite-cusp correction.}
Put \(c=t^{36}\), an actual cusp
loop in \(K_3\), and write
\(u_j=t^j u t^{-j}\). Define
\(b_t=u_0u_1\cdots u_{35}\).
In \(Y\) this is the element
\(N_{36}(t)v\) of \(\mathcal L\).
The action of \(t\) on \(U\)
has order dividing 36. Thus
\(b_{t,Y}\) centralizes \(t\).
Postcomposing \(\varphi\) with
\(\operatorname{Int}(b_t)\)
changes its \(Y\)-action by exactly
that inner automorphism, with the
same full homology matrix. Both
\(b_t\) and every one of its
source factors lie in \(S_{57}\),
so the complete D.57 action,
its conjugators, and \(\mathcal W_h\)
are unchanged.

The last cyclic factor is kept
with its exact source relation:
\(u_{36}=cuc^{-1}\). If
\(U_1=u_1\cdots u_{35}\), direct
multiplication in the class-two
kernel gives
\[
 \begin{aligned}
 b_t\,t b_t^{-1}t^{-1}
   &=uU_1(u_1\cdots u_{35}u_{36})^{-1}\\
   &=[u,U_1]\,[c,u]^{-1}.
 \end{aligned}
 \tag{D.62.17}
\]
Both factors belong to \(\mathcal K\).
The second has a cusp homology
factor, so (D.62.12) kills it under
every selected self or cross
projector. The first has projected
exterior coordinate
\[
 e_\Delta\left(v\wedge\sum_{j=1}^{35}t^jv\right)
       =e_\Delta\left(v\wedge N_{36}(t)v\right).
 \tag{D.62.18}
\]
This proves the second actual
source map in (D.62.7). It does
not replace \(c\) by one.

Its polar expression vanishes on
coinvariants: in the second sum
\([w\wedge t^jv]=[v\wedge t^{-j}w]\),
and the index set \(0,\ldots,35\)
is invariant under negation modulo
the action period. The two sums
therefore cancel. The map vanishes
on \(\ker N_{36}(t)\), and hence
factors through the cyclic norm.
Its full group norm is zero as
well. In the sum over \(G\) and
\(\langle t\rangle\), ordered pairs
of distinct vectors occur in both
orders and diagonal exterior terms
are zero. This proves that both
actual correction images lie in
the asserted norm kernel.

Products of the displayed source
corrections add their central errors.
For the \(s\)-corrections, the
updated involution has the same
conjugation action on the kernel:
its difference is central there.
For the cusp corrections, their
commutators with \(t\) are central
in that kernel. Thus any prescribed
sum of their coinvariant columns
is realized by a finite product
of the actual words, with exactly
the retained-source properties just
proved.

\emph{The full affine marking and the evaluated image.}
The final norm calculation uses
the actual map of D.60.2. If
\(v_*=(1,1,1,0)\) and \(Q_g\)
is the retained four-dimensional
matrix of \(g\in Q\), then
\[
 s=(0,S),\qquad r=(v_*,R),\qquad
 t=(Q_Sv_*,SR).
 \tag{D.62.19}
\]
For the selected character \(\ell\),
let \(g_j\) be the retained pair
coset representatives. If
\(gg_j=g_bh\), the action of
\((a,g)\) on that summand is
the recorded semilinear matrix
of \(h\), followed by the scalar
\[
                 \zeta_3^{\ell Q_{g_b^{-1}}a}.
 \tag{D.62.20}
\]
This is the literal induced
representation of the affine
semidirect product: its character
on the output summand is
\(\ell Q_{g_b^{-1}}\). It proves
the marking used in (D.62.9).

The program
\path{d_self_wedge_cusp_correction.py}
constructs (D.62.20) on every real
basis vector. It verifies
\(s^2=r^3=1\), the actual product
\(t=sr\), all 36 norm powers,
and the displayed affine value
of \(t^{18}\). The coefficient
action has period twelve, but
the program uses the full source
norm of length 36. This retains
the original cyclic source map.

It then checks every polar column
and both cyclic factorizations.
In the recorded ordered real basis
\(v_0,\ldots,v_{n-1}\), independent
columns are as follows. A column
number in this table denotes an
actual basis vector of the
induced homology module, with the
coset representatives and finite
integral reduction fixed in D.61.
\[
 \begin{array}{c|c|c|c}
 \text{self type}&\text{map and vector}
      &\text{coinvariant image}&\text{basis order}\\\hline
 (0,0,0,1),i=4&\eta_t(v_0),\eta_t(v_1)
      &\mathtt{0x4},\mathtt{0xc}&(0,9,16,17)\\
 (0,0,1,1)&\eta_s(v_{96}),\eta_s(v_{100}),\eta_t(v_1)
      &\mathtt{0x4},\mathtt{0x1},\mathtt{0x2}&(0,10,16)\\
 (0,0,1,2)&\eta_s(v_0),\eta_s(v_3),\eta_t(v_0)
      &\mathtt{0x1},\mathtt{0x4},\mathtt{0x2}&(0,9,16)
 \end{array}
 \tag{D.62.21}
\]
The last column gives indices in
the D.61 Hermitian basis. The
hexadecimal entries are the literal
binary coordinates on the resulting
coinvariant basis, and all other
columns are retained in the
certificate. The ordinary norm on
the first row has kernel spanned
by its last two coordinates.
Thus its two displayed columns
are a basis of that kernel.
In each new row the three columns
are a basis of the whole
three-dimensional coinvariant space.
The four copies in each new
orbit have identical literal
matrices, as proved and checked
in D.61. These columns prove the
joint image assertion in every
copy, and give (D.62.8).

This evaluates the correction
map completely on the relevant
kernel. For example, an old
kernel class \(a\mathtt{0x4}
+b\mathtt{0x8}\) is removed by
\((a+b)\eta_t(v_0)+b\eta_t(v_1)\).
For a new row, the displayed
three coordinate basis vectors
give the three coefficients
directly. All coefficients lie
in \(\mathbb F_2\). Taking the
corresponding source factors with
coefficient one constructs the
finite products in (D.62.5) and
(D.62.15). No unproved search
or solvability condition remains
on these twenty-six coordinates.

\emph{The actual boundary and the larger source.}
Before these corrections, D.55
applies to (D.62.14). By
\(eT_{61}=e\), the compact
orientation image on every
selected self block is fixed.
Compactification is injective
on these blocks. Therefore
\[
                   \nu([\Pi_{62}z])=0.
 \tag{D.62.22}
\]
The actual arithmetic boundary
belongs to exactly the kernel
spanned in (D.62.21). Choose
the finite source products with
the coefficients just described.
Their errors cancel that class.

They retain the same homology
matrix. Every previous selected
invariant block still has its
zero boundary by the same norm
and compactification argument.
More explicitly, the new
commutator vectors lie in a
single excluded self factor;
the additional term in (D.62.17)
has a cusp factor killed by all
the selected projectors. The
\(\Delta\)-moving part has zero
coinvariants by odd averaging.
In every retained central
component, the remaining error
representative is thus in the
actual augmentation image
\[
                    (1-s)M+s(1-r)M=I_GM.
 \tag{D.62.23}
\]
Choose central conjugator corrections
with these values. They leave the
action and conjugators on \(Y\),
\(X_*\) and \(\mathcal W_h\)
unchanged. The final cusp error
lies in \(\mathcal J_{62}\).

The relation
\(\Pi_{62}\bigwedge^2\bar T_{61}=\Pi_{62}\)
preserves its kernel, and the old
projectors retain the same property.
Thus the source map factors
through \(\mathcal T_2/\mathcal J_{62}\).
It is an automorphism, since its
kernel Frattini action is the
invertible reduction of \(T_{61}\)
and its finite base action is
identity. The retained full
homology and \(\mathcal W_h\)
action agree with it on their
original marked common quotient.
Their fibre product is exactly
\(X/\mathcal J_{62}\), by the
same source-kernel intersection
as in D.58 and D.61. This proves
(D.62.4) and the explicit action
comparison (D.62.6).

The integral homology matrix
has not changed. The \(Y\)-action
has changed only by the inner
element in (D.62.5), and its
image in \(X_*\) and
\(\mathcal W_h\) is one.
This proves preservation of
the prescribed literal parameter
and all comparison maps claimed
in the statement.

The orthogonal projectors in
(D.62.2) prove all the direct
sums and exact source kernels.
The dimension calculations are
\[
 \begin{aligned}
 \dim\mathcal C_{62}&=512+8\cdot768=6656,\\
 \dim(\mathcal C_{62})_G&=4+8\cdot3=28,\\
 331776-6656&=325120,\qquad562-28=534.
 \end{aligned}
 \tag{D.62.24}
\]
The source image \(R_{50}\to Z\)
is onto. A nonzero vector in
the last specified residual and
one of its actual source preimages
give a nontrivial element of
\(S_{62}/(R_{50}\cap L_2)\).
This proves the stated remaining
kernel without any cofinality claim.

\emph{The exponent-eight identity check and its exact scope.}
The independent program
\path{d_self_wedge_word_identities.py}
uses the full Hall coordinates of
the free class-two groups with
abelian exponent four or eight
and commutator exponent two or
four, respectively. It evaluates
the complete 36-factor cusp word
with the actual cyclic wrap
\(u_{36}=cuc^{-1}\), and the
two-factor source-involution word.
Both (D.62.16) and (D.62.17) hold
with their full signed central
coordinates in both groups.
In particular the cusp term is
retained before its projected
vanishing is used. This checks
the literal square/commutator
identity needed in the proof.
The larger arithmetic
exponent-eight kernel has not
been removed by that universal
word test. The remaining actual
arithmetic boundary is the one
on \((\mathcal J_{62})_G\), with
the exact maps (D.62.10)--(D.62.11).
\end{proof}
